\begin{tcolorbox}[title=Problem 2, breakable]
    Prove Corollary 2.9 and find an example to show that the corollary does not hold without the 
        finiteness condition.
\end{tcolorbox}

\begin{theorem}
    Let $\tau \in \mathcal{L}(V, W)$, where $\dim(V) = \dim(W) < \infty$.
    Then $\tau$ is injective if and only if it is surjective.
\end{theorem}

\begin{tcolorbox}[title=Problem 3, breakable]
    Let $\tau \in \mathcal{L}(V, W)$. Prove that $\tau$ is an isomorphism if and only
    if it carries a basis for $V$ to a basis for $W$.
\end{tcolorbox}

\begin{proof}
    Suppose $\tau$ is an isomorphism.
    Let $\mathcal{B}_V = \{v_i\}$ be a basis of $V$.
    Let $\tau(\mathcal{B}_V) = T$.
    If $T$ is not a linearly independent set of vectors,
        then $\sum c_i \tau(v_i) = 0$ with not all $c_i = 0$, but then applying the inverse
        gives $\sum c_i v_i = 0$, which contradicts the linear independence of $\mathcal{B}_V$.
    Also, for any $w \in W$, write $\tau^{-1}(w) = \sum c_i v_i$, so
        $w = \sum c_i \tau(v_i)$ and $T$ spans $W$.
    Thus $T$ is a basis for $W$.

    Conversely, suppose $\tau$ carries a basis $\mathcal{B}_V = \{v_i\}$ for $V$ to a basis for $W$.
    For any $w \in W$, we have $w = \sum c_i \tau(v_i) = \tau\left(\sum c_i v_i\right)$,
        so $\tau$ is surjective.
    If $\tau(v) = 0$, write $v = \sum c_i v_i$, then $\sum c_i \tau(v_i) = 0$, so all
        $c_i = 0$ by linear independence, and $v = 0$. So $\tau$ is injective.
    Thus $\tau$ is an isomorphism.
\end{proof}

\begin{tcolorbox}[title=Problem 4, breakable]
    If $\tau \in \mathcal{L}(V_1, W_1)$ and $\sigma \in \mathcal{L}(V_2, W_2)$ we define the external direct
        sum $\tau \boxplus \sigma \in \mathcal{L}(V_1 \boxplus V_2, W_1 \boxplus W_2)$ by 
    \[(\tau \boxplus \sigma)((v_1, v_2)) = (\tau v_1, \sigma v_2),\]
    show that $\tau \boxplus \sigma$ is a linear transformation.
\end{tcolorbox}

\begin{proof}
    Let $(v_1, v_2), (v'_1, v'_2) \in V_1 \boxplus V_2$.
    Furthermore, let $a, b \in F$.
    Then we have a vector $a(v_1, v_2) + b(v'_1, v'_2) = (av_1 + bv'_1,\ av_2 + bv'_2) \in V_1 \boxplus V_2$, and
    \begin{align*}
        (\tau \boxplus \sigma)\big(a(v_1, v_2) + b(v'_1, v'_2)\big)
            &= (\tau \boxplus \sigma)\big((av_1 + bv'_1,\ av_2 + bv'_2)\big) \\
            &= \big(\tau(av_1 + bv'_1),\ \sigma(av_2 + bv'_2)\big) \\
            &= \big(a\tau v_1 + b\tau v'_1,\ a\sigma v_2 + b\sigma v'_2\big) \\
            &= a(\tau v_1, \sigma v_2) + b(\tau v'_1, \sigma v'_2) \\
            &= a(\tau \boxplus \sigma)((v_1, v_2)) + b(\tau \boxplus \sigma)((v'_1, v'_2)).
    \end{align*}
    Thus $\tau \boxplus \sigma$ is a linear transformation.
\end{proof}

\begin{tcolorbox}[title=Problem 5, breakable]
    Let $V = S \oplus T$. Prove that $S \oplus T \cong S \boxplus T$.
\end{tcolorbox}

\begin{proof}
    Let $\mathcal{B}_S = \{b_i\}_{i \in I}$ and $\mathcal{B}_T = \{b'_j\}_{j \in J}$ be bases for $S$ and $T$.
    Consider the map $\tau \colon V = S \oplus T \to S \boxplus T$ defined by
    \[
        \tau(v) = (s, t), \quad \text{where } v = s + t \text{ for } s \in S,\ t \in T.
    \]
    First, we want to show that $\tau$ is well-defined. Since $V = S \oplus T$, every
    $v \in V$ can be written uniquely as $v = s + t$ for some $s \in S$ and $t \in T$.
    Thus $\tau(v) = (s, t)$ is well-defined. Since $\mathcal{B}_S$ and $\mathcal{B}_T$ are
    bases, we can write $s$ and $t$ uniquely as
    \[
        s = \sum_{i \in I} f(i)\, b_i, \qquad t = \sum_{j \in J} f'(j)\, b'_j,
    \]
    where $f(i) = 0$ for all but finitely many $i \in I$ and $f'(j) = 0$ for all but finitely many $j \in J$,
    since $s$ and $t$ are finite linear combinations of elements of $\mathcal{B}_S$ and $\mathcal{B}_T$.
    Next, we want to show that $\tau$ is linear. Let $v, w \in V$ and $a \in F$, and suppose
    that $v = s + t$ and $w = s' + t'$ for $s, s' \in S$ and $t, t' \in T$. Then
    \[
        av + w = (as + s') + (at + t'),
    \]
    where $as + s' \in S$ and $at + t' \in T$ since $S$ and $T$ are subspaces. Since this
    decomposition is unique,
    \[
        \tau(av + w) = (as + s',\ at + t') = a(s, t) + (s', t') = a\,\tau(v) + \tau(w).
    \]
    Thus $\tau$ is linear.
    Now, we want to show that $\tau$ is bijective.
    Suppose that $\tau(v) = (0_S, 0_T)$. Then
    \[
        \tau(v) = (s, t) = \left( \sum_{i \in I} f(i)\, b_i,\ \sum_{j \in J} f'(j)\, b'_j \right) = (0_S, 0_T).
    \]
    Then $\sum_{i \in I} f(i)\, b_i = 0_S$ and $\sum_{j \in J} f'(j)\, b'_j = 0_T$.
    Since $\mathcal{B}_S$ and $\mathcal{B}_T$ are linearly independent, $f(i) = 0$ for all
    $i \in I$ and $f'(j) = 0$ for all $j \in J$. Thus $s = 0_S$ and $t = 0_T$, so
    \[
        v = s + t = 0_V.
    \]
    Thus $\ker \tau = \{0_V\}$, and $\tau$ is injective.
    Now, suppose that $(s, t) \in S \boxplus T$. Then $v = s + t \in V$, and since this
    decomposition is unique, $\tau(v) = (s, t)$. Thus $\tau$ is surjective.
    Thus $\tau$ is a linear bijection, and $S \oplus T \cong S \boxplus T$.
\end{proof}

\begin{tcolorbox}[title=Problem 6, breakable]
    Let $V = A + B$ and consider the external direct sum $E = A \boxplus B$.
    Define a map $\tau : A \boxplus B \mapsto V$ by $\tau(v, w) = v + w$.
    Show that $\tau$ is linear. What is the kernel of $\tau$? When is $\tau$ an isomorphism.
\end{tcolorbox}

\begin{proof}
    Let $x, y \in A \boxplus B$ such that $x = a_1 + b_1, y = a_2 + b_2$ and take $c, d \in F$.
    Then $\tau(cx + dy) = cx + dy = (a_1 + b_1) + (a_2 + b_2) = $
\end{proof}